\section{Consequences of the finite-memory conjugacy}\label{sec:measurement}

We now record the consequences of Theorem~\ref{thm:two-sided-conjugacy} for
finite blocks, path-space laws and continuous potentials. These are consequences
of the conjugacy; the arithmetic input has already been isolated in
Section~\ref{sec:stab:conjugacy}.

\subsection{Block presentations above threshold}\label{sec:measurement:sufficiency}

Fix $m\ge 3$ and $n\ge m$.  Write
\[
\mathcal Y_{m,n}:=\Phi_{m,n}\bigl(\{0,1\}^{n+m-1}\bigr)\subset X_m^n
\]
for the legal length-$n$ stabilized output blocks, and let
\[
\Psi_{m,n}:\mathcal Y_{m,n}\to\{0,1\}^{n+m-1}
\]
denote the inverse of $\Phi_{m,n}$ from Theorem~\ref{thm:block-invertibility}.

\begin{theorem}[Equivalence of block presentations above threshold]\label{thm:exact-sufficiency}
For the deterministic equivalence of raw and stabilized block presentations
above threshold, see \cite[Theorem~6.1]{MaZhang2026Nonlinearity}.
\end{theorem}

\subsection{Path-space equivalence}

The same inverse is local on two-sided trajectories. Thus the many-to-one
character of $\Fold_m$ on isolated windows does not lead to a loss of
path-space information for $m\ge3$.

\begin{corollary}[Path-space equivalence and relative entropy rates]\label{cor:path-space-divergence}
For path-space equivalence and divergence-rate preservation above threshold,
see \cite[Corollary~6.2]{MaZhang2026Nonlinearity}.
\end{corollary}

\begin{remark}[Local fibers versus trajectory information]
The entropy identity in Proposition~\ref{thm:kl-ledger} concerns a single
window and records unresolved multiplicity inside the finite fibers of
$\Fold_m$. Corollary~\ref{cor:path-space-divergence} is a different statement:
overlapping windows carry enough consistency information to reconstruct the
whole trajectory above the threshold.
\end{remark}

\subsection{Continuous potentials and finite-range interactions}\label{sec:measurement:renorm}

We finally state the thermodynamic consequence of the conjugacy. This is the
standard transport of potentials through a topological conjugacy, written here
to identify the induced finite-range interaction explicitly.

\begin{theorem}[Transport of continuous potentials under the conjugacy]\label{thm:continuous-potentials}
For transport of continuous potentials under the conjugacy, see
\cite[Theorem~6.4]{MaZhang2026Nonlinearity}.
\end{theorem}

\begin{proposition}[Finite-range energetics under local inversion]\label{prop:finite-range-renorm}
For finite-range energetics under local inversion, see
\cite[Proposition~6.6]{MaZhang2026Nonlinearity}.
\end{proposition}

\begin{remark}[Prototype observables]
Taking $\phi(a_0):=a_0$ gives exact recovery of occupancy sums, and taking
$\phi(a_0,a_1):=a_0a_1$ gives exact recovery of nearest-neighbour pair sums.
Thus the occupancy and nearest-neighbour energy examples used below are
instances of Proposition~\ref{prop:finite-range-renorm}; they do not require
additional thermodynamic assumptions.
\end{remark}
